\documentclass[11pt,a4paper,openany]{ctexbook}
\usepackage[a4paper,margin=1.78cm,headheight=16pt]{geometry}
\usepackage{amsmath,amssymb,mathtools,bm}
\usepackage{booktabs,longtable,tabularx,array,multirow}
\newcolumntype{L}[1]{>{\raggedright\arraybackslash}p{#1}}
\newcolumntype{Y}{>{\raggedright\arraybackslash}X}
\setlength{\extrarowheight}{2pt}
\usepackage{enumitem,tcolorbox,tikz,xcolor,fancyhdr,hyperref,microtype,titlesec,lastpage}
\tcbuselibrary{breakable,skins,theorems}
\usetikzlibrary{arrows.meta,positioning,calc,shapes.geometric,fit,matrix}
\definecolor{ink}{RGB}{25,25,25}
\definecolor{deep}{RGB}{42,58,73}
\definecolor{soft}{RGB}{245,247,249}
\definecolor{line}{RGB}{150,158,166}
\definecolor{accent}{RGB}{104,76,55}
\definecolor{warn}{RGB}{123,76,67}
\definecolor{greenish}{RGB}{57,92,76}
\definecolor{purple}{RGB}{87,72,116}
\hypersetup{colorlinks=true,linkcolor=deep,urlcolor=accent,pdftitle={数学一线性代数心智模型手册 v2}}
\setlength{\parindent}{2em}
\setlength{\parskip}{0.30em}
\setlength{\emergencystretch}{3em}
\renewcommand{\arraystretch}{1.22}
\setlength{\tabcolsep}{3pt}
\setlist{itemsep=0.17em,topsep=0.35em,leftmargin=2em}
\pagestyle{fancy}
\fancyhf{}
\lhead{\small\color{deep}\textbf{数学一 · 线性代数}}
\rhead{\small\color{deep}\leftmark}
\cfoot{\small 第 \thepage\ 页}
\renewcommand{\headrulewidth}{0.4pt}
\titleformat{\chapter}[display]{\normalfont\huge\bfseries\color{deep}}{\filleft\Large 第\thechapter 篇}{0.4ex}{\titlerule\vspace{0.8ex}\filright}
\titleformat{\section}{\Large\bfseries\color{deep}}{\thesection}{0.6em}{}
\titleformat{\subsection}{\large\bfseries\color{ink}}{\thesubsection}{0.6em}{}
\newtcolorbox{corebox}[1]{breakable,colback=soft,colframe=deep,title=\textbf{#1},boxrule=0.8pt,arc=2mm}
\newtcolorbox{methodbox}[1]{breakable,colback=white,colframe=greenish,title=\textbf{#1},boxrule=0.7pt,arc=1.5mm}
\newtcolorbox{warnbox}[1]{breakable,colback=white,colframe=warn,title=\textbf{#1},boxrule=0.7pt,arc=1.5mm}
\newtcolorbox{examBox}[1]{breakable,colback=white,colframe=accent,title=\textbf{数学一训练：#1},boxrule=0.7pt,arc=1.5mm}
\newtcolorbox{boundarybox}[1]{breakable,colback=soft,colframe=line,title=\textbf{概念边界：#1},boxrule=0.6pt,arc=1.5mm}
\newtcolorbox{examplebox}[1]{breakable,colback=white,colframe=purple,title=\textbf{贯穿例：#1},boxrule=0.7pt,arc=1.5mm}
\newcommand{\key}[1]{\textbf{\color{deep}#1}}
\newcommand{\eng}[1]{\textit{#1}}
\let\cleardoublepage\clearpage
\tikzset{
 atlas/.style={draw=deep,rounded corners=2mm,minimum width=2.85cm,minimum height=1cm,align=center,fill=soft,font=\small},
 flow/.style={-{Latex[length=2mm]},thick,draw=deep},
 dashedflow/.style={-{Latex[length=2mm]},densely dashed,thick,draw=purple}
}

\begin{document}
\frontmatter
\begin{titlepage}
\centering\vspace*{0.8cm}
{\Huge\bfseries\color{deep} 线性代数\par}
\vspace{0.3cm}
{\LARGE 心智模型手册 v2\par}
\vspace{0.18cm}
{\small\bfseries\color{warn}历史快照 / Source Only：不再作为当前 Canonical Owner；现行入口为 \texttt{20\_线性代数/README.md}。\par}
\vspace{0.4cm}
{\large 空间 · 映射 · 坐标 · 变换 · 不变量 · 标准化\par}
\vspace{0.8cm}
\begin{tikzpicture}[
  atlas/.style={draw=deep,rounded corners=2.5mm,text width=3.8cm,minimum height=1.15cm,align=center,fill=soft,font=\small,inner sep=3pt,thick},
  flow/.style={-{Latex[length=2.2mm,width=1.5mm]},thick,draw=deep},
  dashedflow/.style={-{Latex[length=2.2mm,width=1.5mm]},thick,dashed,draw=accent}
]
% 第一排：概念与坐标层 (x=0, 5.8, 11.6)
\node[atlas] (sp) at (0,0) {\textbf{Space (向量空间)}\\\scriptsize 独立方向与自由度};
\node[atlas] (map) at (5.8,0) {\textbf{Linear Map (线性映射)}\\\scriptsize $\ker T$ 与 $\operatorname{Im} T$};
\node[atlas] (coord) at (11.6,0) {\textbf{Coordinates (矩阵表示)}\\\scriptsize 选定基后的坐标化};

% 第二排：变换与分类层 (y=-2.2)
\node[atlas] (trans) at (11.6,-2.2) {\textbf{Transform (坐标变换)}\\\scriptsize 等价 / 相似 / 合同};
\node[atlas] (std) at (5.8,-2.2) {\textbf{Canonical Form (标准形)}\\\scriptsize 最简表示与主轴化};
\node[atlas] (inv) at (0,-2.2) {\textbf{Invariants (不变量)}\\\scriptsize 秩 / 谱 / 惯性指数};

% 正向流
\draw[flow] (sp.east) -- node[above,font=\tiny\color{deep}]{抽象化} (map.west);
\draw[flow] (map.east) -- node[above,font=\tiny\color{deep}]{坐标化} (coord.west);
\draw[flow] (coord.south) -- node[right,font=\tiny\color{deep}]{换基/变量代换} (trans.north);

% 底行反向化简流
\draw[flow] (trans.west) -- node[above,font=\tiny\color{deep}]{坐标化简} (std.east);
\draw[flow] (std.west) -- node[above,font=\tiny\color{deep}]{分类不变量} (inv.east);

% 左侧闭环上升线
\draw[dashedflow] (inv.north) -- node[left,font=\tiny\color{accent}]{提取本征特征} (sp.south);
\end{tikzpicture}
\vfill
{\small 矩阵是坐标语言；消元、对角化和化标准形，都是为了在不破坏目标结构的前提下换一种更容易看清问题的表示。\par}
\end{titlepage}

\chapter*{使用说明：不要把线性代数学成“矩阵计算学”}
\addcontentsline{toc}{chapter}{使用说明：不要把线性代数学成“矩阵计算学”}

\begin{corebox}{本册真正要完成什么}
线性代数最容易被学成六块互不相干的计算：行列式、矩阵、向量组、方程组、特征值、二次型。真正稳定的理解不是再造一套口诀，而是建立一套能反复生成这些章节的世界模型：
\[
\boxed{\text{Space}+\text{Linear Map}+\text{Coordinates}+\text{Transform}+\text{Invariants}}.
\]
本册承担五项任务：
\begin{enumerate}
  \item 解释“向量、矩阵、方程组、特征值、二次型”在同一个线性世界中的位置；
  \item 解释每一种计算为什么出现、它实际改变了什么、又保持了什么；
  \item 建立可直接进入题目的概念边界，不只写 $A\neq B$，而说明题目信号与错误后果；
  \item 给出贯穿例，让抽象的空间、映射、坐标、不变量真正落到矩阵计算；
  \item 留出 Practice Layer：以后做题得到的识别规则、起手动作和检查经验可以继续增长，而不破坏稳定理论骨架。
\end{enumerate}
\end{corebox}

\begin{methodbox}{本册的“四层阅读协议”}
以后看到任何线代题，先把内容分成四层：
\[
\boxed{\text{Object}\rightarrow\text{Representation}\rightarrow\text{Allowed Transform}\rightarrow\text{Invariant}}.
\]
这里箭头不是数学运算，而是\key{思考顺序}：
\begin{itemize}
  \item \key{Object：}真正研究的是向量组、线性方程、线性算子还是二次型？
  \item \key{Representation：}当前看到的矩阵、坐标、方程只是这个对象的哪种表示？
  \item \key{Allowed Transform：}为了简化问题，允许做行变换、等价变换、相似变换还是合同变换？
  \item \key{Invariant：}变换以后，哪个量必须不变，因而能控制最终答案？
\end{itemize}
很多错题不是“算错”，而是第一层对象就判断错，随后把本来属于相似的问题拿合同做，把解集问题拿特征值做。
\end{methodbox}

\begin{methodbox}{符号和箭头怎样读}
本册中：
\begin{itemize}
  \item $V\xrightarrow{T}W$ 表示一个真正的线性映射，$T$ 把定义域 $V$ 中的向量送到陪域 $W$；
  \item $A\to B$ 若出现在模型链中，通常表示\key{认知推进或化简方向}，不是等式；
  \item $A\sim B$ 在正文中若明确写“相似”才表示相似矩阵；
  \item $x=Py$ 表示变量/坐标改变，必须继续问 $P$ 是把新坐标变成旧坐标，还是反过来；
  \item “不变量”不是“矩阵数字完全不变”，而是某种变换下代表同一结构的量保持不变。
\end{itemize}
任何公式第一次出现，本册都会配中文解释；公式是压缩后的结论，不承担解释本身。
\end{methodbox}

\tableofcontents
\mainmatter

% ============================================================
\chapter{总图：线性代数到底在研究什么}
% ============================================================

\section{从一个最小世界开始：空间、映射和坐标}
先暂时忘掉行列式和矩阵。设 $V,W$ 是两个线性空间，$T:V\to W$ 满足
\[
T(\alpha u+\beta v)=\alpha T(u)+\beta T(v).
\]
线性意味着：\key{只要知道一组基向量经过 $T$ 后变成什么，就知道整个空间中所有向量怎样被 $T$ 作用。}

若 $V$ 的一组基为 $\mathcal B=(e_1,\dots,e_n)$，$W$ 的一组基为 $\mathcal C=(f_1,\dots,f_m)$，把每个 $T(e_j)$ 在 $\mathcal C$ 下的坐标作为第 $j$ 列，就得到矩阵 $A$。于是抽象关系
\[
V\xrightarrow{T}W
\]
在坐标世界中写成
\[
\boxed{y=Ax}.
\]
因此最关键的第一条边界是：
\[
\boxed{\text{Linear Map}\neq\text{Matrix};\qquad \text{Matrix}=\text{Map + chosen bases 的坐标表示}.}
\]

\section{五个生成模型：六章教材怎样长在同一棵树上}

\subsection{Model 1：空间怎样由独立方向生成}
\[
\boxed{\text{Linear Combination}\to\text{Span}\to\text{Independence}\to\text{Basis}\to\text{Dimension}}
\]
这条链的中文含义是：先问“这些方向能生成什么”；再问“里面有没有冗余”；删掉冗余后得到一组最小坐标骨架；骨架需要多少根，就是空间的自由度。

\subsection{Model 2：线性映射把自由度送到哪里}
\[
\boxed{V\xrightarrow{T}W,\qquad \ker T,\ \operatorname{Im}T,\ \operatorname{rank}T}
\]
输入方向只有两个命运：要么进入 $\ker T$ 被压成零，要么在 $\operatorname{Im}T$ 中留下独立输出方向。因此
\[
\boxed{\dim V=\dim\ker T+\dim\operatorname{Im}T}.
\]
这就是 Rank--Nullity 的真正内容：\key{输入自由度被分成“丢失的自由度”和“保留下来的自由度”。}

\subsection{Model 3：方程组是逆像问题}
\[
Ax=b\iff T(x)=b.
\]
题目不再只是“消元求 $x$”，而是在问：\key{$b$ 是否位于 $T$ 能达到的 image 中？如果能达到，一共有多少种输入会被送到同一个 $b$？}
\[
\boxed{Ax=b\text{ 有解}\iff b\in\operatorname{Im}A}
\]
若 $x_0$ 是一组特解，则所有解为
\[
\boxed{x=x_0+x_h,\qquad x_h\in\ker A}.
\]

\subsection{Model 4：对象不变，坐标表示可以改变}
同一个向量换基后坐标会变；同一个线性算子换基后矩阵会变。线代真正关心的是：
\[
\boxed{\text{Representation changes; structure may stay the same.}}
\]
因此必须区分“对象本身”和“在某组基下写出来的数字”。

\subsection{Model 5：标准化就是寻找最适合对象的坐标}
线代的大量技巧最终都在做同一件事：\key{在允许的变换下，把对象化到最简单的代表形式。}
\begin{itemize}
  \item 一般线性映射：定义域与陪域分别换基，归结到 rank；
  \item 同一空间上的线性算子：同一套基同时用于输入输出，归结到相似与 eigenstructure；
  \item 二次型：变量换基，归结到合同、惯性和正定性。
\end{itemize}

\section{三套标准化不是三种公式，而是三类分类问题}
\begin{longtable}{L{2.4cm}L{3.6cm}L{3.6cm}L{\dimexpr\linewidth-10.17cm\relax}}
\toprule
\textbf{关系} & \textbf{研究对象} & \textbf{坐标变换} & \textbf{核心不变量/标准形}\\
\midrule
等价 & 一般 $V\to W$ 线性映射 & $B=PAQ$，两端可分别换基 & rank；$\begin{psmallmatrix}I_r&0\\0&0\end{psmallmatrix}$\\
相似 & $T:V\to V$ 线性算子 & $B=P^{-1}AP$，同一空间换基 & 特征多项式 / 特征结构；能否对角化\\
合同 & 二次型 $q(x)=x^TAx$ & $B=P^TAP$，同一变量替换进入两个槽位 & rank、正负惯性指数、正定性\\
\bottomrule
\end{longtable}

\begin{corebox}{一个重要的新理解：线代不是“算矩阵”，而是“选择允许的表示变化”}
消元、换基、对角化、配方，看起来是四种技巧，真正不同点在于它们\key{允许改变什么}：
\begin{itemize}
  \item 消元允许把方程作等价重写，目标是保留解集；
  \item 相似允许同一线性算子换基，目标是保留算子结构；
  \item 合同允许二次型换变量，目标是保留二次型的几何类型；
  \item 正交变换比一般可逆变换更严格，还要求保留内积、长度和角度。
\end{itemize}
所以做题的第一步不是“想公式”，而是先问：\key{我当前允许怎样改写对象而不改变题目真正关心的东西？}
\end{corebox}

\section{母模型与教材章节的映射}
\begin{longtable}{L{2.8cm}L{3.9cm}L{4.2cm}L{\dimexpr\linewidth-11.47cm\relax}}
\toprule
\textbf{母模型} & \textbf{教材章节} & \textbf{真正结构} & \textbf{常见计算外壳}\\
\midrule
空间生成 & 向量、向量空间、基、坐标 & span、冗余、自由度 & 相关无关、极大无关组、Schmidt\\
映射自由度 & 矩阵、秩、逆 & kernel/image、信息保留与损失 & 初等变换、rank、inverse\\
逆像问题 & 齐次/非齐次方程组 & 可达性、解集维数、平移结构 & 增广矩阵、基础解系、参数分类\\
坐标改变 & 基变换、相似矩阵 & 同一对象的不同表示 & 过渡矩阵、$P^{-1}AP$\\
标准化与不变量 & 特征值、二次型 & eigenstructure / inertia & 对角化、正交化、配方法、正定判别\\
\bottomrule
\end{longtable}

\section{三条贯穿母例：后面各章会反复回到它们}

\begin{examplebox}{母例 A：rank=2 的三个向量}
令
\[
\alpha_1=(1,1,1)^T,\quad \alpha_2=(1,2,3)^T,\quad \alpha_3=(2,3,4)^T.
\]
因为 $\alpha_3=\alpha_1+\alpha_2$，三个向量只含两个独立方向。以后它会用于解释：span、极大无关组、向量组秩、列空间、线性方程组可达性以及“为什么初等变换是在找独立信息”。
\end{examplebox}

\begin{examplebox}{母例 B：二维平面被压成一条线}
令
\[
A=\begin{pmatrix}1&1\\1&1\end{pmatrix},\qquad T(x,y)=(x+y,x+y).
\]
所有输出都落在 $y=x$ 上，所以 $\operatorname{rank}A=1$；而 $x+y=0$ 的方向被压成零，所以
\[
\ker A=\operatorname{span}\{(1,-1)^T\}.
\]
两个输入自由度中一个被保留、一个被消灭：$2=1+1$。以后它会贯穿 rank、方程组、可逆性与 determinant。
\end{examplebox}

\begin{examplebox}{母例 C：实对称矩阵把特征结构和二次型接起来}
令
\[
S=\begin{pmatrix}2&1\\1&2\end{pmatrix}.
\]
它有两个正交特征方向 $(1,1)^T,(1,-1)^T$，特征值分别为 $3,1$。在线性算子视角，它可以正交相似对角化；在二次型
\[
q(x)=x^TSx
\]
视角，同一个正交坐标变换又能消除交叉项，并立即看出 $q$ 正定。这个例子会解释为什么实对称矩阵是“相似理论”和“合同理论”的交汇点。
\end{examplebox}

% ============================================================
\chapter{向量与空间：先把“自由方向”和“坐标骨架”建立起来}
% ============================================================

\section{线性组合：线性代数最原始的动作}
给定 $v_1,\dots,v_k$，表达式
\[
a_1v_1+\cdots+a_kv_k
\]
不是一个公式技巧，而是在问：\key{允许沿这些方向任意叠加以后，能够到达哪些向量？}
所有可达向量构成
\[
\operatorname{span}\{v_1,\dots,v_k\}.
\]
因此“一个向量能否由另一组向量表示”与“它是否落在这组向量生成的 span 中”是同一件事。

\section{线性相关：本质是表示系统存在冗余}
若
\[
x_1v_1+\cdots+x_kv_k=0
\]
存在非零解，则至少一个方向可以由其他方向线性表示。所谓“相关”不是说向量长得像，而是说\key{生成能力有重复}。

\begin{boundarybox}{线性相关 $\neq$ 两向量成比例}
两个非零向量时，“相关”等价于成比例；但三个及以上向量时，相关只要求某个向量能由其余若干向量组合得到，不要求任意两向量成比例。题目给三维空间中的四个向量时，即使任意两向量都不成比例，也必然相关，因为空间只有三个自由方向。
\end{boundarybox}

\section{极大线性无关组：从一组生成器中抽出信息骨架}
“极大”不是向量数量在全空间里最大，而是：\key{在当前给定向量组内部，已经无法再加入剩余向量而保持无关。}
任何极大无关组都生成与原向量组相同的 span，因此其向量个数相同，这个共同个数就是向量组的秩。

\begin{examplebox}{母例 A：从三个向量中抽骨架}
把 $A=(\alpha_1,\alpha_2,\alpha_3)$ 作行化简：
\[
\begin{pmatrix}1&1&2\\1&2&3\\1&3&4\end{pmatrix}
\to
\begin{pmatrix}1&0&1\\0&1&1\\0&0&0\end{pmatrix}.
\]
主元列来自原矩阵第 1、2 列，所以 $\{\alpha_1,\alpha_2\}$ 是一个极大无关组，向量组秩为 2；第 3 列依赖关系提示 $\alpha_3=\alpha_1+\alpha_2$。

这里有两个不同动作：行变换用于\key{暴露列之间的依赖关系}；真正选极大无关组时，应回到\key{原矩阵的主元列}，不能直接把化简后矩阵的列向量当成原向量组。
\end{examplebox}

\section{向量组等价：不是元素相同，而是生成能力相同}
两个向量组 $A=(\alpha_1,\dots,\alpha_m)$ 与 $B=(\beta_1,\dots,\beta_s)$ 等价，真正含义是
\[
\operatorname{span}(A)=\operatorname{span}(B).
\]
因此常见判据可以理解为：
\begin{itemize}
  \item $A$ 能由 $B$ 表示：$\operatorname{span}(A)\subseteq\operatorname{span}(B)$；
  \item $B$ 也能由 $A$ 表示：反向包含；
  \item 两边同时成立：生成空间相同。
\end{itemize}

\begin{examBox}{向量组等价题的第一步}
题面出现“甲向量组可由乙向量组线性表示”“两向量组等价”“秩是否相同”时，先翻译成 span 包含关系，不要把它误当成坐标逐个相等。若 $A=BC$，只能先得出 $\operatorname{span}(A)\subseteq\operatorname{span}(B)$ 和 $r(A)\le r(B)$；要得到等价，还需要反向表示或秩等附加条件。
\end{examBox}

\section{向量空间与子空间：先看闭合性，再谈基与维数}
线性空间提供“加法”和“数乘”的封闭环境。数学一常见对象多是 $\mathbb R^n$ 的子空间，例如齐次方程解空间、某组向量的 span。

判断集合 $U\subseteq V$ 是否为子空间，最实用的统一判据是：
\[
\forall u,v\in U,\ \forall \alpha,\beta\in\mathbb R,
\qquad \alpha u+\beta v\in U.
\]
它一次性包含加法与数乘封闭，也自动包含零向量。

\begin{boundarybox}{齐次约束天然给子空间，非齐次约束通常不给}
$Ax=0$ 的解集对线性组合封闭，因此是子空间；$Ax=b\ (b\neq0)$ 的解集一般不经过原点，是一个子空间的平移，所以不是线性子空间。这也是为什么“基础解系”只属于齐次方程。
\end{boundarybox}

\section{基与维数：最小无冗余坐标骨架}
一组基必须同时满足：
\[
\boxed{\text{Basis}=\text{Spanning}+\text{Independent}}.
\]
它既能生成整个空间，又没有冗余。维数则是任意一组基所含向量数，代表空间的独立自由度数量。

在 $n$ 维空间中，有一条极实用的“二选一升级规则”：
\begin{itemize}
  \item $n$ 个向量若线性无关，则自动构成基，不必再单独证明能生成；
  \item $n$ 个向量若能生成整个空间，则自动线性无关，不必再单独证明无冗余。
\end{itemize}

\section{坐标不是向量：同一个对象可以有很多数字表示}
设基 $\mathcal B=(\alpha_1,\dots,\alpha_n)$，若
\[
v=x_1\alpha_1+\cdots+x_n\alpha_n,
\]
则 $X=(x_1,\dots,x_n)^T$ 是 $v$ 在基 $\mathcal B$ 下的坐标，记作 $[v]_{\mathcal B}=X$。

同一个 $v$ 换基后，坐标会改变；但向量本身没有改变。这个边界是后面相似矩阵的基础。

\section{过渡矩阵：先固定“谁由谁表示”，再写公式}
设旧基 $\mathcal B=(\alpha_1,\dots,\alpha_n)$，新基 $\mathcal B'=(\beta_1,\dots,\beta_n)$，若
\[
\mathcal B'=\mathcal BP,
\]
则 $P$ 的第 $j$ 列就是新基向量 $\beta_j$ 在旧基 $\mathcal B$ 下的坐标。对同一个向量 $v$，有
\[
v=\mathcal B[v]_{\mathcal B}=\mathcal B'[v]_{\mathcal B'}=\mathcal BP[v]_{\mathcal B'},
\]
所以
\[
\boxed{[v]_{\mathcal B}=P[v]_{\mathcal B'},\qquad [v]_{\mathcal B'}=P^{-1}[v]_{\mathcal B}.}
\]

\begin{methodbox}{基变换题的“方向防错法”}
不要先背 $P$ 还是 $P^{-1}$。先写一句中文：\key{“新基由旧基怎样表示？”} 若 $\mathcal B'=\mathcal BP$，那 $P$ 把新坐标送回旧坐标；因此旧坐标 $=P\times$ 新坐标。这个方向一旦明确，逆矩阵自然出现。
\end{methodbox}

\section{内积：给线性空间增加欧氏几何}
仅有线性结构时，我们只能谈 span、独立、维数；引入内积后，才有长度、夹角、垂直、投影。
\[
\lVert x\rVert=\sqrt{\langle x,x\rangle},\qquad x\perp y\iff \langle x,y\rangle=0.
\]
正交向量组中的非零向量必然线性无关，因为不同方向之间没有“投影泄漏”。

\section{Schmidt 正交化：不是为了复杂公式，而是为了换到干净坐标}
给线性无关组 $\alpha_1,\dots,\alpha_n$，Gram--Schmidt 逐步去掉新向量在已有方向上的投影：
\[
\beta_1=\alpha_1,
\]
\[
\beta_2=\alpha_2-\frac{\langle\alpha_2,\beta_1\rangle}{\langle\beta_1,\beta_1\rangle}\beta_1,
\]
后续同理。最后单位化得到规范正交基。

其本质是：\key{保持原来的 span 不变，但把坐标方向改造成彼此正交、长度为 1 的基。}

\section{正交矩阵：保留几何的换坐标}
若 $Q^TQ=I$，则 $Q^{-1}=Q^T$。更重要的是
\[
\lVert Qx\rVert=\lVert x\rVert,\qquad \langle Qx,Qy\rangle=\langle x,y\rangle.
\]
所以正交变换不改变长度和夹角。其列向量（也等价地行向量）构成规范正交组，且 $\lvert\det Q\rvert=1$。

\begin{boundarybox}{线性无关 $\neq$ 正交 $\neq$ 规范正交}
线性无关只要求“没有冗余”；正交进一步要求两两内积为 0；规范正交还要求每个向量长度为 1。题目给“正交矩阵”时可以立刻用 $Q^{-1}=Q^T$，但普通可逆矩阵不能这样做。
\end{boundarybox}

% ============================================================
\chapter{矩阵与行列式：坐标化之后，怎样合法计算}
% ============================================================

\section{同一个矩阵可以有四种阅读方式}
面对 $m\times n$ 矩阵
\[
A=[a_1\ \cdots\ a_n],
\]
至少有四种常用视角：
\begin{enumerate}
  \item \key{列向量组：}$Ax=x_1a_1+\cdots+x_na_n$，列空间回答哪些输出 $b$ 能被表示；
  \item \key{线性映射：}$T:\mathbb R^n\to\mathbb R^m$，kernel/image 描述输入自由度怎样损失和保留；
  \item \key{方程组系数：}$Ax=b$ 是寻找 $b$ 的所有逆像；
  \item \key{对称方阵时的二次型系数：}$x^TAx$ 描述一个二次量，此时研究对象已经不是线性映射本身。
\end{enumerate}
同一矩阵可以承载不同对象；\key{题目研究哪一个对象，决定你允许用什么理论。}

\section{矩阵乘法：复合线性作用的坐标语言}
若 $B$ 先作用、$A$ 后作用，则
\[
x\xmapsto{B}Bx\xmapsto{A}A(Bx)=(AB)x.
\]
所以矩阵乘法不是任意规定，而是“复合映射”在坐标中的表达。$AB$ 一般不等于 $BA$，也正因为“先做谁后做谁”通常不可交换。

\begin{boundarybox}{$AB=0$ 不能推出 $A=0$ 或 $B=0$}
矩阵没有数域中的零因子性质。$AB=0$ 只说明 $B$ 的 image 全部被 $A$ 的 kernel 吸收：
\[
\operatorname{Im}B\subseteq\ker A.
\]
这条空间解释比死记反例更强。题目一旦出现 $AB=0$，应立即考虑 rank 与 image/kernel 的包含关系。
\end{boundarybox}

\section{方阵的幂：重复作用同一个线性算子}
$A^k$ 表示同一算子连续作用 $k$ 次。直接连乘适合低阶简单矩阵；若 $A$ 可对角化，后面会用
\[
A=P\Lambda P^{-1}\quad\Longrightarrow\quad A^k=P\Lambda^kP^{-1}
\]
把重复复杂作用变成各特征方向上的幂次缩放。

\section{特殊矩阵不是名词表，而是结构提示}
\begin{longtable}{L{2.7cm}L{5.0cm}L{\dimexpr\linewidth-8.27cm\relax}}
\toprule
\textbf{结构} & \textbf{立即可用的性质} & \textbf{题目中意味着什么}\\
\midrule
对角矩阵 & 运算、幂、行列式、特征值都逐元素 & 已经在“最简单坐标”中\\
上/下三角 & 行列式为对角元乘积，特征值也是对角元 & 常用来快速读 det/eigenvalue\\
对称矩阵 & 实特征值，可正交对角化 & eigen 与 quadratic form 的桥梁\\
反对称矩阵 & $A^T=-A$，实数域对角元为 0 & 常与转置运算、二次项消失结合\\
正交矩阵 & $Q^{-1}=Q^T$，保长度角度 & 最安全的几何换基\\
\bottomrule
\end{longtable}

\section{初等行变换：本质是对方程作可逆重写}
三类初等行变换分别对应：交换方程顺序、某方程乘非零常数、某方程加另一方程倍数。它们都不改变方程组的解集。

每次行变换等价于左乘一个可逆初等矩阵 $E$：
\[
A\to EA.
\]
连续行变换就是
\[
A\to E_k\cdots E_1A.
\]
这解释了为什么初等矩阵可逆，也解释了高斯消元为何不会凭空制造或丢失解。

\section{列变换为什么不能在解方程时随便做}
右乘初等矩阵 $Q$ 对应列变换：
\[
A\to AQ.
\]
这相当于重新组合未知变量对应的列方向。如果只对系数矩阵做列变换而不同时跟踪未知量变换，原方程组的 $x$ 已经换了含义。因此：
\begin{itemize}
  \item 求 rank 时，行列初等变换都可用，因为 rank 是等价不变量；
  \item 求 $Ax=b$ 的原变量解时，通常只对增广矩阵作行变换；
  \item 化矩阵等价标准形时，行列变换都可以，因为目标是 rank 分类而不是保留原坐标中的解向量。
\end{itemize}

\section{逆矩阵：可逆意味着信息没有丢失}
若存在 $A^{-1}$ 使
\[
A^{-1}A=AA^{-1}=I,
\]
则 $A$ 的作用可以完全撤销。几何上就是：\key{没有两个不同输入被压到同一个输出，也没有目标方向无法到达。}

\section{可逆矩阵定理：十几个条件其实只说一件事}
对 $n\times n$ 方阵 $A$，下列典型条件等价：
\[
\det A\neq0
\iff r(A)=n
\iff \ker A=\{0\}
\iff \operatorname{Im}A=\mathbb R^n
\]
\[
\iff Ax=0\text{ 只有零解}
\iff Ax=b\text{ 对任意 }b\text{ 唯一可解}
\iff A\text{ 可逆}.
\]
它们不是需要分开背的七条定理，而是同一个结构事实：
\[
\boxed{\text{A loses no degree of freedom.}}
\]

\begin{methodbox}{选择题遇到“可逆的等价命题”怎样做}
不要逐条回忆。先锁定核心：方阵可逆 $\Leftrightarrow$ rank 满 $\Leftrightarrow$ kernel 为零 $\Leftrightarrow$ 每个 $b$ 都唯一可达。再把选项翻译进这个链条。比如“列向量组线性无关”就是 $r(A)=n$；“0 是特征值”则意味着不可逆。
\end{methodbox}

\section{行列式：方阵线性映射的有向体积缩放}
行列式只对方阵定义。它的深层几何意义是：$A$ 把单位平行多面体的有向体积放大为 $\det A$ 倍。于是
\[
\det A=0
\]
意味着 $n$ 维体积被压扁到低维空间，即 rank 小于 $n$。

因此 determinant 与可逆性联系自然出现，而不是人为凑出的条件。

\section{行列式性质：每条都应知道在“体积”上做了什么}
\begin{itemize}
  \item 交换两行：改变坐标方向定向，det 变号；
  \item 某行乘 $k$：沿一个方向缩放 $k$ 倍，det 乘 $k$；
  \item 某行加另一行倍数：剪切不改体积，det 不变；
  \item 三角矩阵：体积缩放由各坐标轴独立缩放相乘，故 det 为对角元乘积；
  \item $\det(AB)=\det A\det B$：复合映射的体积缩放倍数相乘；
  \item $\det A^T=\det A$：行/列两种描述给同一体积缩放。
\end{itemize}

\section{按行列展开：不是必然首选，而是利用稀疏结构}
Laplace 展开把高阶 determinant 降成低阶 determinant。真正适合的场景是某行/列有大量零，或者能通过初等变换先制造零。大行列式不应默认直接展开。

\section{伴随矩阵：把余子式结构组织成求逆工具}
伴随矩阵 $A^*$ 满足
\[
AA^*=A^*A=(\det A)I.
\]
若 $\det A\neq0$，立即得到
\[
A^{-1}=\frac{1}{\det A}A^*.
\]
但伴随矩阵不只是求逆公式，数学一常考它与 rank 的关系。对 $n\times n$ 矩阵（$n\ge2$）：
\[
r(A^*)=
\begin{cases}
 n,&r(A)=n,\\
 1,&r(A)=n-1,\\
 0,&r(A)\le n-2.
\end{cases}
\]
原因是：若满秩，$A^*=\lvert A\rvert A^{-1}$ 当然满秩；若 rank 恰为 $n-1$，至少一个 $(n-1)$ 阶子式非零而 $AA^*=0$，故 $A^*$ 非零且其列都落在一维 kernel 中；若 rank 更低，所有 $(n-1)$ 阶子式都为 0，所以 $A^*=0$。

\begin{examBox}{伴随矩阵题不要只背 rank 三段式}
先看 $AA^*=\lvert A\rvert I$。若 $\lvert A\rvert\neq0$，直接化成逆矩阵问题；若 $\lvert A\rvert=0$，再看 $A^*$ 的元素本质上是 $(n-1)$ 阶余子式。这样 rank 结论是从结构推出，而不是孤立口诀。
\end{examBox}

\section{分块矩阵：利用结构，而不是把大矩阵拆着玩}
分块的价值是保留整体乘法结构，同时让零块、对角块、三角块发挥作用。常用原则：
\begin{itemize}
  \item 分块加法和数乘要求对应块同型；
  \item 分块乘法仍遵守“行块乘列块”，但块尺寸必须匹配；
  \item 块对角矩阵的行列式、逆、幂可逐块处理；
  \item 看到重复块、零块或单位块时，先问能否把大问题降成小问题。
\end{itemize}

\begin{boundarybox}{矩阵可逆 $\neq$ 每个元素非零}
可逆性看的是整体 rank 是否满，而不是元素是否为零。单位矩阵有大量零仍可逆；全 1 矩阵没有零却 rank 为 1（高阶时不可逆）。元素级直觉在矩阵世界里非常危险。
\end{boundarybox}

% ============================================================
\chapter{秩：独立信息维数，连接向量、矩阵和方程组}
% ============================================================

\section{rank 的六种视角是同一个量}
对矩阵 $A$，rank 可以从多个层面理解：
\begin{itemize}
  \item 列向量组的秩；
  \item 行向量组的秩；
  \item 列空间 $\operatorname{Col}(A)$ 的维数；
  \item 线性映射 image 的维数；
  \item 最大非零子式的阶数；
  \item 阶梯形中主元个数。
\end{itemize}
这些不是六个独立定义，而是“独立信息有几维”在几何、映射、代数、计算层面的不同观测方式。

\section{为什么行秩等于列秩}
行初等变换不会改变行空间的维数，同时把矩阵化成阶梯形；阶梯形中的非零行数就是行秩，也是主元列数量。原矩阵主元列给出列空间的一组骨架，因此行秩与列秩都等于主元个数。

在数学一里不必把证明推到更抽象层，但必须留下这个认知：\key{rank 是矩阵内部唯一的一份“独立维数”，不会因为你从行看还是从列看而改变。}

\section{Rank--Nullity：自由度怎样守恒}
对 $m\times n$ 矩阵 $A$，映射 $T:\mathbb R^n\to\mathbb R^m$ 满足
\[
\boxed{n=r(A)+\dim\ker A}.
\]
这条公式的做题意义非常直接：
\begin{itemize}
  \item 只要知道 rank，就知道齐次方程自由变量个数；
  \item 只要知道基础解系有几个向量，就能反推 rank；
  \item 参数改变 rank 时，解空间维数会同步跳变。
\end{itemize}

\begin{examplebox}{母例 B：二维压成一维}
\[
A=\begin{pmatrix}1&1\\1&1\end{pmatrix}
\]
的列空间由 $(1,1)^T$ 生成，所以 rank=1；kernel 是 $x+y=0$，维数 1。于是
\[
2=1+1.
\]
这个等式不是算术巧合，而是整个二维输入的每个自由方向最终只有“保留”或“被压掉”两种命运。
\end{examplebox}

\section{初等变换为什么保持 rank}
行变换左乘可逆矩阵 $P$，列变换右乘可逆矩阵 $Q$：
\[
A\mapsto PAQ.
\]
可逆变换本身不会丢失自由度，所以
\[
\boxed{r(PAQ)=r(A)}.
\]
这也是矩阵等价关系用 rank 分类的根本原因。

\section{rank 运算规律：不要只背不等式，要看 image 被怎样限制}
常用结论包括：
\[
r(A+B)\le r(A)+r(B),
\]
\[
r(AB)\le \min\{r(A),r(B)\}.
\]
第二条的空间解释最重要：
\[
\operatorname{Im}(AB)=A(\operatorname{Im}B),
\]
所以输出维数不可能超过 $B$ 先产生的维数，也不可能超过 $A$ 自己最大能保留的维数。

若 $P,Q$ 可逆，则
\[
r(PA)=r(AQ)=r(A),
\]
因为左右乘可逆矩阵只是在输入或输出端换坐标。

\section{\texorpdfstring{$AB=0$}{AB=0} 时的 rank 约束}
由
\[
\operatorname{Im}B\subseteq\ker A
\]
可得
\[
r(B)\le \dim\ker A=n-r(A)
\]
（尺寸匹配时），因此常见结论
\[
\boxed{r(A)+r(B)\le n.}
\]
这种题的最快入口不是“对 $AB=0$ 展开元素”，而是把它翻译成 image 被 kernel 吸收。

\section{rank 与 determinant：一个是维数，一个是方阵满秩探针}
行列式能回答方阵“是不是满秩”，但不能告诉你 rank 到底是 $n-1$ 还是更低；rank 则对任意矩阵都定义，直接度量独立维数。

\begin{boundarybox}{非零子式判 rank 的真正含义}
$r(A)=r$ 表示：至少存在一个 $r$ 阶子式非零，说明确实能找到 $r$ 个独立行列方向；而所有 $r+1$ 阶子式都为零，说明不可能再增加一个独立方向。含参数 rank 题常利用这个“某个子式何时变零”来定位 rank 跳变点。
\end{boundarybox}

\section{含参数 rank 题：先找“结构跳变点”}
对含参数 $a$ 的矩阵，rank 通常在绝大多数 $a$ 下保持一个常值，只在某些使主元消失或关键子式为零的参数上下降。推荐流程：
\begin{enumerate}
  \item 先做不需要除以含参数表达式的消元；
  \item 找出哪些步骤会出现 $a-c$ 这样的关键因子；
  \item 分 $a\neq c$ 与 $a=c$ 两类；
  \item 每一类重新确认主元数量，而不是把一般情形的除法硬套到特殊点。
\end{enumerate}

% ============================================================
\chapter{线性方程组：把“求解”升级为“研究逆像结构”}
% ============================================================

\section{齐次方程 \texorpdfstring{$Ax=0$}{Ax=0}：直接在求 kernel}
齐次方程解集
\[
N(A)=\{x:Ax=0\}=\ker A
\]
天然是子空间。若 $A$ 有 $n$ 个未知量、rank 为 $r$，则
\[
\dim N(A)=n-r.
\]
基础解系就是这个 kernel 的一组基，因此必含 $n-r$ 个线性无关解向量。

\section{齐次方程有非零解的真正条件}
\[
Ax=0\text{ 有非零解}\iff r(A)<n.
\]
若 $A$ 是 $n\times n$ 方阵，还等价于
\[
\det A=0.
\]
注意：方阵时 determinant 能快速判，但一般 $m\times n$ 方程组仍应使用 rank。

\section{非齐次方程 \texorpdfstring{$Ax=b$}{Ax=b}：先问 \texorpdfstring{$b$}{b} 是否可达}
列向量视角下
\[
Ax=x_1a_1+\cdots+x_na_n=b,
\]
所以有解等价于
\[
b\in\operatorname{Col}(A).
\]
计算上等价于
\[
\boxed{r(A)=r(A\mid b)}.
\]
若增广列带来了新主元，说明 $b$ 不在原列空间中，于是无解。

\section{非齐次通解：特解 + 看不见的方向}
若 $Ax_0=b$，又有 $Ax_h=0$，则
\[
A(x_0+x_h)=b.
\]
反过来任意两个非齐次解之差都属于 kernel。因此
\[
\boxed{\mathcal S_b=x_0+\ker A}.
\]
几何上它是一个与 kernel 平行的仿射集合。

\begin{boundarybox}{“有无穷多解”不是因为方程多或未知数多}
真正条件是：先相容，即 $r(A)=r(A\mid b)$；再看 $r(A)<n$，是否存在自由变量。只有“相容 + kernel 非零”才产生无穷多解。
\end{boundarybox}

\section{基础解系怎样构造：自由变量就是 kernel 的坐标}
把阶梯形中的自由变量依次设为标准单位组合，就得到 $n-r$ 个基础解向量。这不是机械技巧：每个自由变量代表 kernel 中一个独立可移动方向，基础解系就是这些自由方向的一组坐标骨架。

\section{含参数方程组：先分类 rank，再谈具体解}
参数方程题的真正对象不是“把未知数解出来”，而是研究
\[
r(A),\qquad r(A\mid b),\qquad n-r(A)
\]
怎样随参数变化。

\begin{methodbox}{参数方程组四步协议}
\begin{enumerate}
  \item 先把参数视为符号，做安全消元；
  \item 找 rank 可能下降的参数点；
  \item 对每个参数区间/特殊点比较 $r(A)$ 与 $r(A\mid b)$；
  \item 只有确定“无解/唯一/无穷多”以后，再在需要的分支求具体通解。
\end{enumerate}
这比一开始追着 $x_1,x_2,x_3$ 解稳定得多。
\end{methodbox}

\section{Cramer 法则：满秩方阵的显式坐标公式}
当 $A$ 为 $n\times n$ 且 $\det A\neq0$ 时，方程 $Ax=b$ 有唯一解，可用 Cramer 法则
\[
x_i=\frac{D_i}{D}.
\]
它的理论意义是“可逆 $\Rightarrow$ 任意 $b$ 唯一可达”；但含参数、欠定、超定、解结构题中，rank 与行变换更通用。

\section{矩阵方程 \texorpdfstring{$AX=B$}{AX=B}：其实是一次解多个右端项}
若
\[
AX=B=[b_1\ \cdots\ b_k],
\]
并写 $X=[x_1\ \cdots\ x_k]$，则等价于同时求
\[
Ax_j=b_j,\qquad j=1,\dots,k.
\]
若 $A$ 可逆，当然 $X=A^{-1}B$；若不可逆，则每一列 $b_j$ 都必须落在 $\operatorname{Col}(A)$ 中才能相容。

\section{\texorpdfstring{$XA=B$ 与 $AX=B$}{XA=B 与 AX=B} 不可随便对称处理}
矩阵乘法不交换。$XA=B$ 若 $A$ 可逆，应右乘 $A^{-1}$ 得 $X=BA^{-1}$；$AX=B$ 则左乘得 $X=A^{-1}B$。考试中“逆矩阵乘哪边”错误，本质上是忘了矩阵代表有方向的复合作用。

\section{两个方程组解集关系：用 kernel/span 语言统一}
若两个齐次方程组为 $Ax=0$、$Bx=0$，那么“所有 A 的解也是 B 的解”就是
\[
\ker A\subseteq\ker B.
\]
在具体计算题中，常转化为：$B$ 的每个约束能由 $A$ 的约束推出，或者通过解空间基检验包含关系。这个视角比把两组方程逐条对照更稳定。

\begin{examBox}{方程组题的结果校验}
得到答案后至少检查三件事：
\begin{enumerate}
  \item 基础解系向量数量是否等于 $n-r(A)$；
  \item 每个基础解向量代回是否满足 $Ax=0$；
  \item 非齐次通解中的特解是否满足 $Ax_0=b$，齐次部分是否真的被 $A$ 压成 0。
\end{enumerate}
如果数量先不对，后面表达再漂亮也一定有错。
\end{examBox}

% ============================================================
\chapter{坐标改变：等价、相似、合同为什么必须严格分开}
% ============================================================

\section{一般线性映射为什么对应矩阵等价}
设 $T:V\to W$。如果在定义域 $V$ 换一组基，同时在陪域 $W$ 也独立换一组基，那么矩阵会发生左右两侧不同的可逆变换，形式可写为
\[
B=PAQ
\]
（具体 $P,Q$ 方向取决于基变换约定，但本质是“两端可分别换基”）。因此一般线性映射的分类只剩一个核心维数：rank。

通过合适基，可以把 rank 为 $r$ 的映射写成最简单形：
\[
\boxed{
\begin{pmatrix}
I_r&0\\0&0
\end{pmatrix}}.
\]
它直接说出：$r$ 个方向被保留，其余输入方向被压掉。

\section{为什么线性算子对应相似，而不是一般等价}
若 $T:V\to V$，定义域和陪域是同一个空间。我们不是随意在输入端用一套基、输出端再用另一套无关基，而是用\key{同一套新基}描述“变换前”和“变换后”的向量。

设旧基 $\mathcal B$，新基 $\mathcal B'=\mathcal BP$。若 $A,B$ 分别是 $T$ 在两基下的矩阵，则
\[
\boxed{B=P^{-1}AP}.
\]
所以相似的真正含义是：
\[
\boxed{\text{same operator, different coordinates}.}
\]

\section{相似公式完整推导：不要背 \texorpdfstring{$P^{-1}$}{P inverse} 在哪边}
由
\[
T(\mathcal B)=\mathcal BA,\qquad T(\mathcal B')=\mathcal B'B,\qquad \mathcal B'=\mathcal BP,
\]
有
\[
T(\mathcal B')=T(\mathcal BP)=T(\mathcal B)P=\mathcal BAP,
\]
另一方面
\[
T(\mathcal B')=\mathcal B'B=\mathcal BPB.
\]
所以 $AP=PB$，即
\[
\boxed{B=P^{-1}AP}.
\]
这同时解释了对角化时为什么 $P$ 的列应该放特征向量：因为 $AP=P\Lambda$ 正是在说每个新基向量都是 eigenvector。

\section{为什么二次型对应合同}
二次型
\[
q(x)=x^TAx
\]
不是线性映射 $x\mapsto Ax$，而是把一个向量代入两次后得到标量。若变量换成
\[
x=Py,
\]
则
\[
q(Py)=y^TP^TAPy.
\]
所以新矩阵是
\[
\boxed{B=P^TAP}.
\]
$P^T$ 的出现不是“另一套记忆公式”，而是同一个变量 $x=Py$ 同时进入了二次型的左右两个槽位。

\begin{center}
\begin{tikzpicture}[
  relnode/.style={draw=deep,rounded corners=2.5mm,text width=3.8cm,minimum height=1.15cm,align=center,fill=soft,font=\small,inner sep=3pt,thick},
  arrow/.style={-{Latex[length=2.2mm,width=1.5mm]},thick,draw=deep},
  dasharrow/.style={-{Latex[length=2.2mm,width=1.5mm]},thick,dashed,draw=accent}
]
% 第一行：三大研究对象 (x=0, 5.8, 11.6)
\node[relnode] (obj_map) at (0,0) {\textbf{线性映射 $T: V\to W$}\\\scriptsize 异基转换 (两端独立换基)};
\node[relnode] (obj_op) at (5.8,0) {\textbf{线性算子 $T: V\to V$}\\\scriptsize 同基转换 (同一套基)};
\node[relnode] (obj_quad) at (11.6,0) {\textbf{二次型 $q(x)=x^TAx$}\\\scriptsize 变量代换 (槽位同代换)};

% 第二行：三大矩阵变换关系 (y=-2.3)
\node[relnode] (rel_eq) at (0,-2.3) {\textbf{矩阵等价}\\\scriptsize $B = PAQ$ ($P,Q$ 可逆)};
\node[relnode] (rel_sim) at (5.8,-2.3) {\textbf{矩阵相似}\\\scriptsize $B = P^{-1}AP$ ($P$ 可逆)};
\node[relnode] (rel_cong) at (11.6,-2.3) {\textbf{矩阵合同}\\\scriptsize $B = P^TAP$ ($P$ 可逆)};

% 第三行：核心分类不变量 (y=-4.6)
\node[relnode] (inv_rank) at (0,-4.6) {\textbf{秩 (Rank)}\\\scriptsize 独立输出维数 $r(A)$};
\node[relnode] (inv_spec) at (5.8,-4.6) {\textbf{谱 (Spectrum)}\\\scriptsize 特征值 / 迹 / 行列式};
\node[relnode] (inv_iner) at (11.6,-4.6) {\textbf{惯性指数 (Inertia)}\\\scriptsize 正负平方项个数 $(p,q)$};

% 关系连线
\draw[arrow] (obj_map.south) -- node[left,font=\tiny\color{deep}]{分类动作} (rel_eq.north);
\draw[arrow] (obj_op.south) -- node[left,font=\tiny\color{deep}]{分类动作} (rel_sim.north);
\draw[arrow] (obj_quad.south) -- node[left,font=\tiny\color{deep}]{分类动作} (rel_cong.north);

\draw[arrow] (rel_eq.south) -- node[left,font=\tiny\color{deep}]{唯一分类不变量} (inv_rank.north);
\draw[arrow] (rel_sim.south) -- node[left,font=\tiny\color{deep}]{保持本征结构} (inv_spec.north);
\draw[arrow] (rel_cong.south) -- node[left,font=\tiny\color{deep}]{保持符号惯性} (inv_iner.north);

% 正交变换汇合线 (底部正交路由)
\draw[dasharrow] (rel_sim.south) -- ++(0,-0.8) -| node[below,pos=0.25,font=\tiny\color{accent}]{正交变换 $Q^TAQ = Q^{-1}AQ$ (相似与合同在正交基下完美统一)} (rel_cong.south);
\end{tikzpicture}
\end{center}

\section{三种关系究竟分别保持什么}
\begin{longtable}{L{2.2cm}L{3.7cm}L{3.4cm}L{\dimexpr\linewidth-9.97cm\relax}}
\toprule
\textbf{关系} & \textbf{一定保持} & \textbf{不一定保持} & \textbf{题面最强信号}\\
\midrule
等价 & rank & eigenvalues、det（非同型时甚至没意义） & 初等行列变换、秩、标准形\\
相似 & rank、det、trace、characteristic polynomial、eigenvalues & 元素、特征向量坐标本身 & 特征值、对角化、同一线性变换换基\\
合同（实对称） & rank、正负惯性指数、正定性 & 一般特征值、trace、det 的具体值 & 二次型、配方法、正交变换、正定\\
\bottomrule
\end{longtable}

\begin{warnbox}{一个高频误区：合同矩阵的特征值一般不保持}
正交合同 $Q^TAQ$ 因为 $Q^{-1}=Q^T$，此时恰好也是相似，所以特征值保持；但一般可逆 $P$ 的合同 $P^TAP$ 不是相似变换，特征值不需要保持。二次型真正保持的是惯性，不是 eigenvalue 的具体数值。
\end{warnbox}

% ============================================================
\chapter{特征值与特征向量：找到算子自己的方向}
% ============================================================

\section{特征向量：不被“转向”的方向}
若
\[
Av=\lambda v,\qquad v\neq0,
\]
那么 $A$ 作用在 $v$ 上只改变长度和可能的方向正负，不把它转向别的独立方向。因此：
\begin{itemize}
  \item eigenvector 是算子的不变方向；
  \item eigenvalue 是沿该方向的伸缩倍数；
  \item $\lambda=0$ 意味着这个 eigenvector 被压进 kernel，因此 $A$ 不可逆。
\end{itemize}

\section{为什么要解 \texorpdfstring{$\det(\lambda I-A)=0$}{det(lambda I - A) = 0}}
条件
\[
(A-\lambda I)v=0
\]
要有非零解，必须
\[
r(A-\lambda I)<n,
\]
对方阵等价于
\[
\boxed{\det(\lambda I-A)=0}.
\]
所以特征方程本质上是在寻找：\key{哪些缩放因子 $\lambda$ 会让 $A-\lambda I$ 突然失去可逆性，从而出现不变方向。}

\section{特征值的常用结构性质}
若 $A$ 为 $n\times n$，则：
\begin{itemize}
  \item 特征值之和（按代数重数）等于 $\operatorname{tr}A$；
  \item 特征值之积等于 $\det A$；
  \item $A$ 与 $A^T$ 特征值相同；
  \item 若 $A$ 可逆，则 $A^{-1}$ 的特征值为 $1/\lambda$；
  \item 若 $Av=\lambda v$，则 $A^kv=\lambda^kv$，更一般地多项式 $p(A)v=p(\lambda)v$。
\end{itemize}
这些不是互不相关的公式：它们都来自“在 eigenvector 方向上，矩阵运算退化成标量运算”。

\section{不同特征值的特征向量为什么自动无关}
若 $v_i$ 对应互不相同的 $\lambda_i$，则这些 $v_i$ 线性无关。直觉上：算子在这些方向上的缩放规则不同，不可能存在非平凡线性关系还能在作用前后同时成立。

这条性质给出一个重要充分条件：\key{$n$ 阶矩阵若有 $n$ 个互不相同特征值，则一定可对角化。}

\section{代数重数与几何重数：重根不等于不能对角化}
对特征值 $\lambda$：
\begin{itemize}
  \item \key{代数重数} $m_a$：$\lambda$ 作为特征多项式根的重数；
  \item \key{几何重数} $m_g$：特征子空间 $\ker(A-\lambda I)$ 的维数。
\end{itemize}
总有
\[
1\le m_g\le m_a.
\]
矩阵可对角化当且仅当所有特征值提供的独立特征方向总数达到 $n$，等价于每个特征值都有
\[
m_g=m_a.
\]

\begin{examplebox}{同样是重特征值，命运完全不同}
\[
A_1=\begin{pmatrix}2&0\\0&2\end{pmatrix},\qquad
A_2=\begin{pmatrix}2&1\\0&2\end{pmatrix}.
\]
两者特征值都只有 $2$，代数重数都是 2。$A_1-2I=0$，eigenspace 维数为 2，所以可对角化；$A_2-2I$ 的 kernel 只有一维，因此只有一个独立特征方向，不可对角化。

所以题目看到“重根”，下一步必须求 $r(A-\lambda I)$ 或 eigenspace 维数，不能立即下结论。
\end{examplebox}

\section{对角化：找到与算子最适配的坐标轴}
若存在 $n$ 个线性无关特征向量 $v_1,\dots,v_n$，令
\[
P=(v_1,\dots,v_n),\qquad \Lambda=\operatorname{diag}(\lambda_1,\dots,\lambda_n),
\]
则
\[
AP=P\Lambda,
\]
从而
\[
\boxed{P^{-1}AP=\Lambda}.
\]
在这组 eigenbasis 中，复杂算子被拆成各独立方向上的纯缩放。

\section{对角化最常见的用途：矩阵幂和矩阵函数}
若 $A=P\Lambda P^{-1}$，则
\[
A^k=P\Lambda^kP^{-1}.
\]
这把重复矩阵乘法变成对角元的幂。做题时看到“求 $A^n$”“递推由矩阵控制”，如果 eigenstructure 简单，应优先考虑对角化。

\section{实对称矩阵为什么是数学一最好的矩阵}
若 $A^T=A$，则：
\begin{enumerate}
  \item 特征值全为实数；
  \item 不同特征值对应的特征向量正交；
  \item 每个重特征值的特征子空间内部可以 Schmidt 正交化；
  \item 最终存在正交矩阵 $Q$ 使
  \[
  \boxed{Q^TAQ=\Lambda}.
  \]
\end{enumerate}
因为 $Q^{-1}=Q^T$，这里既是相似对角化，也是合同对角化。这正是特征值理论与二次型理论真正交汇的地方。

\begin{examplebox}{母例 C：$S=\begin{psmallmatrix}2&1\\1&2\end{psmallmatrix}$}
特征值为 $3,1$，对应方向 $(1,1)^T,(1,-1)^T$。单位化后令
\[
Q=\frac1{\sqrt2}\begin{pmatrix}1&1\\1&-1\end{pmatrix},
\]
则
\[
Q^TSQ=\begin{pmatrix}3&0\\0&1\end{pmatrix}.
\]
在线性算子视角：新坐标轴正好沿着两个不变方向；在二次型视角：交叉耦合被完全消除，而且两个主轴方向上的“强度”分别为 3 和 1。
\end{examplebox}

\section{Jordan 标准形放在哪里：结构扩展，不挤占数学一主线}
若特征向量不够，矩阵不能完全拆成独立缩放方向。更高等的线性代数会用广义特征向量把它化成 Jordan 形，即“特征值缩放 + 最简单链式耦合”。

数学一不要求 Jordan 计算；这里只保留一个心智作用：\key{“不可对角化”并不意味着完全没有结构，而是说明算子不能被拆成彼此完全独立的 eigen-directions。}

% ============================================================
\chapter{二次型：从交叉耦合到主轴、惯性与正定性}
% ============================================================

\section{二次型不是线性映射：研究对象已经换了}
一般二次型可写为
\[
q(x)=\sum_{i=1}^n a_{ii}x_i^2+2\sum_{i<j}a_{ij}x_ix_j=x^TAx,
\]
其中总可以取 $A=A^T$。

这里要特别注意：若题面交叉项写成 $6x_1x_2$，矩阵对称位置应填 $a_{12}=a_{21}=3$，因为 $x^TAx$ 会把两边贡献加起来。

\begin{boundarybox}{二次型矩阵的非对角元要“除以 2”}
题面 $c_{ij}x_ix_j$（$i\neq j$）若没有显式写成 $2a_{ij}x_ix_j$，构造对称矩阵时要令 $a_{ij}=a_{ji}=c_{ij}/2$。这是二次型最常见的低级失分点之一。
\end{boundarybox}

\section{二次型的母问题：怎样去掉坐标之间的耦合}
交叉项 $x_ix_j$ 表示不同坐标方向耦合。化标准形就是找新变量 $x=Py$，使
\[
q(x)=y^TDy=d_1y_1^2+\cdots+d_ny_n^2.
\]
因此标准化的核心不是“把式子写漂亮”，而是\key{寻找主轴，让不同方向的二次贡献彼此独立。}

\section{合同变换完整推导}
若 $x=Py$，则
\[
q(x)=(Py)^TA(Py)=y^T(P^TAP)y.
\]
所以新矩阵
\[
\boxed{B=P^TAP}.
\]
若 $P$ 可逆，则 $A,B$ 合同，代表同一个二次型在不同变量坐标下的表示。

\section{标准形与规范形：先去耦合，再只保留符号类型}
标准形是
\[
d_1y_1^2+\cdots+d_ry_r^2
\]
（其余系数可为 0）；进一步对非零变量做缩放，可化为规范形
\[
z_1^2+\cdots+z_p^2-z_{p+1}^2-\cdots-z_{p+q}^2.
\]
标准形保留具体非零系数；规范形只保留正、负、零三类数量。

\section{惯性定理：合同分类真正不变的是正负方向数量}
任意可逆实合同变换下，标准形中的：
\begin{itemize}
  \item 正平方项个数 $p$；
  \item 负平方项个数 $q$；
  \item 零平方项个数 $n-p-q$
\end{itemize}
保持不变。于是 rank 为 $p+q$。

这给出二次型最深的分类：\key{不同坐标下系数可以变化，但“有多少正方向、多少负方向、多少退化方向”不变。}

\section{配方法：代数上直接寻找能消掉交叉项的新变量}
例如
\[
q(x_1,x_2)=x_1^2+4x_1x_2+5x_2^2
\]
可写成
\[
q=(x_1+2x_2)^2+x_2^2.
\]
令
\[
y_1=x_1+2x_2,\qquad y_2=x_2,
\]
即可得到标准形 $y_1^2+y_2^2$。

配方法的本质是：\key{用代数方式构造一个可逆变量替换，使交叉耦合消失。} 它不要求先求特征值，常适合系数结构简单的题。

\section{正交变换：用实对称矩阵的 eigenbasis 找主轴}
因为二次型矩阵可取实对称 $A$，存在正交矩阵 $Q$ 使
\[
Q^TAQ=\Lambda.
\]
令 $x=Qy$，则
\[
q=\lambda_1y_1^2+\cdots+\lambda_ny_n^2.
\]
这个标准形的系数直接就是 $A$ 的特征值。

\begin{methodbox}{配方法 vs 正交变换：不是“两个做法谁更高级”}
\begin{itemize}
  \item 配方法：只要求找一个可逆变量替换，计算可能更短，但新坐标一般不保持长度和角度；
  \item 正交变换：要求新基规范正交，几何意义最清楚，标准形系数直接是特征值；
  \item 若题目明确要求“正交变换”，必须求正交特征向量并单位化；只配平方不能替代要求；
  \item 若只要求“化标准形”，应根据系数结构选择成本最低的方法。
\end{itemize}
\end{methodbox}

\section{正定：所有非零方向上的二次量都为正}
实对称矩阵 $A$ 正定，定义为
\[
\boxed{x^TAx>0\qquad \forall x\neq0.}
\]
这不是“矩阵元素都是正数”，而是对\key{所有方向}的整体要求。

常见等价判据：
\begin{itemize}
  \item 所有特征值都 $>0$；
  \item 正惯性指数 $p=n$；
  \item 所有顺序主子式 $\Delta_1,\dots,\Delta_n$ 全正（Sylvester 判据）；
  \item 二次型可经可逆变换化成全正平方和。
\end{itemize}

\section{为什么顺序主子式能判正定}
Sylvester 判据看似是一组 determinant 口诀，实际是在逐级检查前 $k$ 个坐标方向张成的子空间上，二次型是否始终保持正。所有层级都通过，最终整个空间上没有负方向或零方向。

\section{正定与特征值：把“所有方向为正”压缩到主轴方向}
对实对称 $A$，正交对角化后
\[
q=\lambda_1y_1^2+\cdots+\lambda_ny_n^2.
\]
由于正交变换覆盖所有非零方向，$q$ 对所有非零 $y$ 为正，当且仅当每个 $\lambda_i>0$。这就是为什么正定判据会自然与 eigenvalue 连接。

\begin{examplebox}{母例 C 再看一次：同一个 $S$ 的两种身份}
\[
S=\begin{pmatrix}2&1\\1&2\end{pmatrix}
\]
的 eigenvalues 是 $3,1$，都正，所以 $S$ 正定。相应二次型
\[
2x_1^2+2x_1x_2+2x_2^2
\]
经正交主轴变换变成
\[
3y_1^2+y_2^2.
\]
这说明“正定”不是观察原式里系数正不正，而是在真正主轴方向上，所有二次强度是否都为正。
\end{examplebox}

\section{二次型与高等数学 Hessian 的接口}
多元函数在驻点附近有二阶近似
\[
f(x_0+h)\approx f(x_0)+\frac12h^THh,
\]
其中 $H$ 是 Hessian。于是多元极值的二阶判别，本质上要判断二次型 $h^THh$ 的正负性：
\begin{itemize}
  \item $H$ 正定：驻点附近各方向都向上，局部极小；
  \item $H$ 负定：各方向都向下，局部极大；
  \item $H$ 不定：存在上升和下降方向，鞍点。
\end{itemize}
这里高数调用线代的“二次型正定性”，不需要在高数手册重复整个合同理论。

% ============================================================
\chapter{全书贯通：一个矩阵如何在不同章节中换身份}
% ============================================================

\section{矩阵不是一个问题，矩阵是多个问题的共同表示载体}
给一个方阵 $A$，题目可能从完全不同的角度问：
\begin{itemize}
  \item 它的列向量能生成多大空间？——向量组/rank；
  \item $Ax=b$ 是否有解？——image/preimage；
  \item $Ax=0$ 有多少自由方向？——kernel/nullity；
  \item 作为 $T:V\to V$ 的表示，有哪些不变方向？——eigenstructure；
  \item 作为 $q(x)=x^TAx$ 的系数，几何类型是什么？——congruence/inertia。
\end{itemize}
所以“看到矩阵先算 determinant”是没有对象意识的表现。

\section{贯通链一：rank 是最强的中间接口}
rank 同时连接：
\[
\boxed{\text{vector dependence}}
\]
\[
\boxed{\text{image dimension}}
\]
\[
\boxed{\text{nullity}=n-r}
\]
\[
\boxed{\text{equation consistency/solution dimension}}
\]
\[
\boxed{\text{invertibility when square}}
\]
因此线代题一旦出现“相关性、方程解数、参数分类、可逆性”，rank 往往是最值得先检查的控制量。

\section{贯通链二：可逆性把 determinant、rank、方程和 eigenvalue 接起来}
对 $n\times n$ 方阵：
\[
\boxed{
A^{-1}\text{ exists}
\iff \det A\neq0
\iff r(A)=n
\iff \ker A=\{0\}
\iff 0\text{ 不是特征值}
}
\]
\[
\boxed{
\iff Ax=b\text{ 对任意 }b\text{ 唯一可解}
\iff \text{列向量组构成 }\mathbb R^n\text{ 的一组基}.
}
\]
这是一条特别值得在选择题里反复调用的“等价命题高速公路”。

\section{贯通链三：实对称矩阵连接正交、特征值和二次型}
实对称矩阵 $A$ 同时拥有：
\[
A^T=A
\Rightarrow
\text{real eigenvalues}
\Rightarrow
\text{orthogonal eigenbasis}
\Rightarrow
Q^TAQ=\Lambda.
\]
因此它既能解决线性算子的对角化，又能解决二次型的主轴化；若进一步所有 $\lambda_i>0$，又直接得到正定。

\section{贯通链四：标准化不是为了“算完”，而是为了暴露结构}
\begin{itemize}
  \item 阶梯形暴露主元、自由变量、rank；
  \item 对角矩阵暴露独立 eigen-directions 与重复作用；
  \item 二次型标准形暴露正、负、零方向；
  \item 规范正交基暴露几何上干净的坐标轴。
\end{itemize}
所以线代的很多算法都可视为：
\[
\boxed{\text{Complex Representation}\to\text{Canonical Representation}\to\text{Read Structure Directly}.}
\]

% ============================================================
\chapter{概念边界：每一个“不等于”都要能改变你的做题动作}
% ============================================================

\begin{longtable}{L{1.7cm}@{}>{\centering\arraybackslash}p{0.35cm}@{}L{1.7cm}L{3.3cm}L{3.3cm}L{\dimexpr\linewidth-10.72cm\relax}}
\toprule
\textbf{概念 A} & & \textbf{概念 B} & \textbf{真正区别} & \textbf{题目信号} & \textbf{混淆后常见错误}\\
\midrule
向量 & $\neq$ & 坐标 & 抽象对象 vs 某组基下的数字表示 & “在基 $\alpha$ 下坐标”“过渡矩阵” & 把换基后的坐标变化当成向量变化\\
矩阵 & $\neq$ & 线性映射 & 表示 vs 抽象作用 & “同一线性变换在不同基” & 看到矩阵不同就认为对象不同\\
线性相关 & $\neq$ & 两两成比例 & 多向量表示有冗余 vs 两向量特殊情形 & 三个以上向量组 & 错判相关性\\
极大无关组 & $\neq$ & 任意无关组 & 前者还要能生成原向量组 span & “求极大无关组” & 少选方向导致 span 变小\\
向量组等价 & $\neq$ & 向量逐个相同 & 生成空间相同 & “互相线性表示” & 忽略 span 关系\\
行变换 & $\neq$ & 列变换 & 左乘可逆矩阵 vs 右乘可逆矩阵 & 解方程/求 rank & 求原变量解时误改未知量含义\\
rank & $\neq$ & det & 独立维数 vs 方阵体积缩放标量 & 非方阵、自由度 vs 满秩 & 对非方阵谈 determinant\\
列空间 & $\neq$ & kernel & 可达到的输出 vs 被压成 0 的输入 & $Ax=b$ / $Ax=0$ & 把 $b$ 放到错误空间\\
有解 & $\neq$ & 唯一解 & 相容性 vs kernel 是否为零 & $r(A)=r(A\mid b)$ 与 $r(A)=n$ & 只判相容就说唯一\\
齐次解空间 & $\neq$ & 非齐次解集 & 子空间 vs 子空间平移 & “基础解系”“特解+齐次解” & 给非齐次解集谈基\\
可逆 & $\neq$ & 元素都非零 & 整体 rank 满 vs 元素级性质 & determinant/rank & 用元素符号判断可逆\\
等价/相似 & $\neq$ & 合同 & 研究对象和允许换基方式不同 & rank/eigen/quadratic form & 不变量乱套\\
代数重数 & $\neq$ & 几何重数 & 多项式根次数 vs eigenspace 维数 & 重特征值 & 错判对角化\\
重特征值 & $\neq$ & 不可对角化 & 重根可能仍有足够特征向量 & 求 $r(A-\lambda I)$ & 看到重根直接否定对角化\\
相似 & $\neq$ & 相等 & 同一算子不同基表示 vs 数字逐项相同 & $P^{-1}AP$ & 把相似当等号可随意代元素\\
正交 & $\neq$ & 线性无关 & 正交是更强的几何条件 & Schmidt/内积 & 把一般基当规范正交基\\
合同 & $\neq$ & 相似 & 二次型换变量 vs 算子换基 & $P^TAP$ / $P^{-1}AP$ & 错认为合同保持 eigenvalue\\
标准形 & $\neq$ & 规范形 & 对角系数任意非零实数 vs 只保留 $\pm1,0$ & 二次型 & 少做/多做变量缩放\\
正定 & $\neq$ & 元素全正 & 所有非零方向上 $x^TAx>0$ & 特征值/主子式 & 看矩阵元素符号判正定\\
\bottomrule
\end{longtable}

\section{边界不是背诵表，而是“题目分流器”}
例如题目写“$A,B$ 有相同特征值”，不能反推出相似，因为 eigenvalue 只是相似不变量的一部分；题目写“$A,B$ 合同且 $A$ 正定”，则可以立即推出 $B$ 正定，因为正定性属于合同不变量。每个“不等于”真正的价值，是让你知道\key{什么信息足以推出什么，什么信息绝对还不够。}

% ============================================================
\chapter{统一做题控制：把线代从“算式反应”变成“结构决策”}
% ============================================================

\begin{corebox}{线代十问}
\begin{enumerate}
  \item \key{Object：}现在研究的是向量组、矩阵、方程解集、线性算子还是二次型？
  \item \key{Space：}对象在哪个空间里？定义域/陪域维数是多少？
  \item \key{Goal：}题目要的是生成能力、自由度、可达性、解集、特征结构还是正负性？
  \item \key{Representation：}当前数字是哪个基下的坐标？矩阵是哪种对象的表示？
  \item \key{Control Quantity：}rank、det、nullity、eigenvalue、inertia 中谁最直接控制目标？
  \item \key{Allowed Transform：}允许行变换、列变换、等价、相似、合同还是正交变换？
  \item \key{Canonical Form：}最值得化成阶梯形、对角形、标准形还是规范正交基？
  \item \key{Execute：}现在才进行消元、求 eigen、配方、Schmidt 等计算。
  \item \key{Interpret：}算出的 rank/eigenvector/solution basis 到底在空间中意味着什么？
  \item \key{Verify：}维数、数量、代回、不变量是否与结构一致？
\end{enumerate}
\end{corebox}

\section{向量组题：先问“生成什么”和“有多少冗余”}
题面出现“线性表示、相关无关、极大无关组、向量组秩、等价”时，统一入口：
\[
\boxed{\text{Build columns}\to\text{Find rank/pivots}\to\text{Interpret span and redundancy}.}
\]
不要一看到“相关”就机械算 determinant；非方阵或向量数量不等于空间维数时，rank 更通用。

\section{矩阵可逆/伴随题：先走可逆等价链}
先问 $\lvert A\rvert$、rank、kernel、0 特征值之间有什么直接信息。若出现 $A^*$，优先使用
\[
AA^*=\lvert A\rvert I
\]
和余子式阶数，不要把 $A^*$ 当普通矩阵重新从头算。

\section{含参数方程组：先分类，不要一上来求未知数}
\[
\boxed{\text{safe elimination}\to\text{rank breakpoints}\to r(A)\text{ vs }r(A\mid b)\to\text{solution type}\to\text{specific solution}}.
\]
这是最值得长期固化成操作习惯的一类题。

\section{特征值题：先判断“求值、判对角化、求幂”是哪一种}
\begin{itemize}
  \item 只求 eigenvalue：优先利用三角结构、trace、det、已知特征关系；
  \item 判可对角化：重根时必须检查 eigenspace 维数；
  \item 求 $A^n$：若可对角化，优先 $P\Lambda^nP^{-1}$；
  \item 实对称：直接想到正交对角化，不要浪费时间纠结可对角化条件。
\end{itemize}

\section{二次型题：先判断目标是“标准形”还是“正定性”}
若只判正定，未必需要完整求特征向量；可以根据结构选顺序主子式、特征值或配方。若要求正交变换化标准形，则必须进入实对称矩阵的 eigenbasis 路线。

\section{选择题证明题：先找“不变量”通常比直接计算快}
例如：
\begin{itemize}
  \item 相似矩阵的 det/trace/eigenvalue 一定相同；
  \item 合同矩阵的正负惯性指数一定相同；
  \item 初等变换/等价矩阵的 rank 相同；
  \item 可逆左/右乘不改变 rank；
  \item $AB=0$ 立即转成 $\operatorname{Im}B\subseteq\ker A$。
\end{itemize}
很多抽象选择题真正考的是“你知不知道当前关系下什么不能变”。

\section{线代统一检查协议}
\begin{methodbox}{完成计算后至少做这六项检查}
\begin{enumerate}
  \item \key{维数：}kernel 基向量个数是否为 $n-r$？
  \item \key{代回：}特征向量是否满足 $Av=\lambda v$？基础解是否满足 $Ax=0$？
  \item \key{数量：}对角化时是否真的有 $n$ 个独立 eigenvectors？
  \item \key{不变量：}相似前后 trace/det/eigenvalue 是否一致？合同前后 inertia 是否一致？
  \item \key{正交：}若题目要求正交矩阵，列向量是否单位化且两两正交？
  \item \key{正定：}结论是否针对所有非零方向，而不是只看几个元素？
\end{enumerate}
\end{methodbox}

% ============================================================
\appendix
\chapter{2026 数学一线性代数完整覆盖清单}
% ============================================================

\section{行列式}
\begin{itemize}
  \item 行列式的概念与基本性质；
  \item 行列式按行（列）展开定理；
  \item 利用性质、初等变换、三角结构和展开计算行列式。
\end{itemize}

\section{矩阵}
\begin{itemize}
  \item 矩阵概念及单位、数量、对角、三角、对称、反对称矩阵；
  \item 线性运算、乘法、转置、方阵幂、方阵乘积行列式；
  \item 逆矩阵、可逆充要条件、伴随矩阵；
  \item 初等变换、初等矩阵、矩阵等价；
  \item 矩阵的秩及用初等变换求 rank、inverse；
  \item 分块矩阵及其运算。
\end{itemize}

\section{向量}
\begin{itemize}
  \item 向量线性组合与线性表示；
  \item 线性相关、线性无关；
  \item 极大线性无关组、向量组秩及与矩阵秩的关系；
  \item 向量空间、子空间、基、维数、坐标；
  \item 基变换与坐标变换；
  \item 内积、Schmidt 正交化、规范正交基、正交矩阵。
\end{itemize}

\section{线性方程组}
\begin{itemize}
  \item Cramer 法则；
  \item 齐次方程非零解充要条件；
  \item 非齐次方程有解充要条件；
  \item 齐次方程基础解系、通解、解空间；
  \item 非齐次方程解的性质、结构与通解；
  \item 用初等行变换求解线性方程组。
\end{itemize}

\section{特征值与特征向量}
\begin{itemize}
  \item 特征值、特征向量概念、性质与计算；
  \item 相似变换、相似矩阵及性质；
  \item 可相似对角化的充要条件与求法；
  \item 实对称矩阵的特征值、特征向量性质与正交对角化。
\end{itemize}

\section{二次型}
\begin{itemize}
  \item 二次型及矩阵表示；
  \item 合同变换、合同矩阵；
  \item 二次型的 rank、惯性定理；
  \item 标准形、规范形；
  \item 正交变换与配方法化标准形；
  \item 正定二次型、正定矩阵及判别法。
\end{itemize}

\begin{methodbox}{覆盖基准与边界}
范围按 2026 数学（一）考试大纲公开转载校验。Jordan 标准形、最小多项式、一般双线性型、四个基本子空间的正交补理论等只在本册中作为理解扩展或不展开，不计入数学一核心主线。正式复习仍应以当年官方考试大纲与题设模型为准。
\end{methodbox}

% ============================================================
\chapter{一页压缩：线性代数最终应该留下什么}
% ============================================================

\begin{corebox}{第一层：四个真正对象}
\textbf{Space：}线性组合决定能到哪里；无关性判断有没有冗余；基是最小无冗余坐标骨架；维数是自由度数量。

\textbf{Map：}$\ker T$ 告诉你哪些输入方向被压成零；$\operatorname{Im}T$ 告诉你哪些输出方向能够到达；rank 是保留下来的独立输出维数。

\textbf{Coordinates：}矩阵和坐标只是选定基以后的数字表示；换基可以改变表示而不改变抽象对象。

\textbf{Invariant：}真正稳定的判断来自“在允许变换下什么不能变”。
\end{corebox}

\begin{corebox}{第二层：三条最值得记住的结构链}
\[
\boxed{\text{span}\to\text{independence}\to\text{basis}\to\text{dimension}}
\]
\[
\boxed{n=\dim\ker A+r(A)}
\]
\[
\boxed{Ax=b\text{ 的解集}=x_0+\ker A}
\]
前两条描述空间和映射的自由度，第三条把方程组直接接回空间结构。
\end{corebox}

\begin{corebox}{第三层：三类标准化}
\textbf{一般映射：}等价变换，核心看 rank；

\textbf{线性算子：}相似变换，核心看 eigenstructure，最好情况是对角化；

\textbf{二次型：}合同变换，核心看 inertia/definiteness，目标是消除交叉耦合。
\end{corebox}

\begin{corebox}{最后的做题原则}
看到一道线代题，不要先问“公式是什么”，而要问：
\[
\boxed{\text{对象是谁？允许怎样改表示？哪个不变量直接控制答案？}}
\]
消元、求秩、求特征值、配方只是最后的执行层。真正稳定的能力，是能在“空间--映射--坐标--变换--不变量”之间来回切换。
\end{corebox}

\end{document}
